\section{Objet et portée de l'annexe}
\label{annexe-H-objet}

Cette annexe établit le théorème de réduction au contenu survivant,
qui constitue le noyau démonstratif de la mécanique des seuils du
Tome III. L'idée centrale est la suivante : lorsque les variables
rapides et les horloges intermédiaires se sont exactement éteintes avant
la sonde, l'histoire antérieure ne peut plus agir sur la réponse que par
le contenu mnésique qui lui survit.

Le résultat ne postule ni l'existence d'un seuil unique, ni la monotonie
de la réponse, ni celle de l'écriture produite par le préfixe. Il affirme
seulement une factorisation de la réponse par l'état survivant. Les
équations de seuil usuelles, les translations de cartes et la
classification des motifs non monotones apparaissent ensuite comme des
conséquences supplémentaires, valables sous des hypothèses explicites.

Dans tout l'appendice, le couple d'observation $(\mathfrak B,\rho)$ est
déclaré. La classe $\mathfrak B$ peut être minimale : aucune stabilité
par continuation n'est utilisée. La résolution $\rho$ est supposée
compatible avec la sonde considérée au sens de l'Appendice~G.

\begin{encadre}[title={Résultat directeur}]
Dans le régime d'extinction exacte, la réponse à une sonde $\mathcal P$
ne dépend de l'histoire passée qu'à travers le contenu survivant
$\mathbf s_{\mathrm r}$. Il existe donc une application
\[
\Omega_{\mathcal P,\rho}
\]
telle que
\[
\mathcal R_X^{\rho}(\mathcal P)
=
\Omega_{\mathcal P,\rho}(\mathbf s_{\mathrm r}).
\]
La fonction $\Omega_{\mathcal P,\rho}$ est la carte de réponse de la
sonde nue sur l'espace des contenus survivants.
\end{encadre}

\section{État de départ de la sonde}
\label{annexe-H-etat-sonde}

\subsection{Décomposition de l'état hybride}
\label{annexe-H-decomposition}

On écrit l'état complet sous la forme
\[
X=(y,p,\chi,\theta,\mathbf s),
\]
où :
\begin{itemize}
\item $y$ rassemble les variables mécaniques ou géométriques remises à
  une graine canonique avant la sonde ;
\item $p$ désigne la pression ou la charge résiduelle ;
\item $\chi$ désigne la variable d'interface rapide ;
\item $\theta$ désigne l'horloge ou la phase lente intermédiaire ;
\item $\mathbf s$ désigne le contenu mnésique survivant.
\end{itemize}

Cette écriture est fonctionnelle. Elle ne prétend pas que ces variables
constituent l'unique paramétrisation possible de l'état ; elle isole les
composantes dont le statut diffère dans la réduction. Les trois
variables $(p,\chi,\theta)$ doivent être éteintes, tandis que
$\mathbf s$ est précisément la composante autorisée à survivre.

Pour une famille alignée dans une direction unitaire $\mathbf u_1$, on
écrit
\[
\mathbf s=s\,\mathbf u_1,
\qquad s\in I\subseteq\mathbb R.
\]
La réduction devient alors scalaire.

\subsection{La regraine comme opération du protocole}
\label{annexe-H-regraine}

La regraine appartient à l'histoire expérimentale. Elle n'est pas une
modification silencieuse de l'état initial.

\begin{encadre}[title={A-H1 --- Appartenance de la regraine au protocole}]
Toute regraine utilisée pour comparer deux histoires est un épisode
déclaré du protocole. Son action sur chaque composante de l'état est
spécifiée, et elle est concaténée à l'histoire avant la sonde.
\end{encadre}

On note $\mathcal G$ l'épisode de regraine et l'on suppose qu'il impose
une configuration canonique $y_0$ aux variables mécaniques :
\begin{equation}
\label{eq:annexe-H-regraine}
\Phi_{\mathcal G}(y,p,\chi,\theta,\mathbf s)
=
(y_0,p,\chi,\theta,\mathbf s).
\end{equation}
La regraine ne doit donc pas effacer le contenu que l'on cherche à
sonder. Dans le modèle considéré, elle réinitialise la position et la
vitesse de la graine, mais conserve les variables internes héritées.

Cette convention est indispensable. Sans elle, deux expériences
portant le même nom de sonde pourraient partir de configurations
mécaniques différentes, et le contenu survivant ne serait pas un résumé
suffisant de leur état de départ.

\subsection{Régime d'extinction exacte}
\label{annexe-H-extinction}

On appelle \emph{régime d'extinction exacte} l'ensemble des états qui,
au moment de la regraine, vérifient
\begin{equation}
\label{eq:annexe-H-extinction}
 p=0,
 \qquad
 \chi=0,
 \qquad
 \theta=0.
\end{equation}
On le note
\[
\Sigma_{\mathrm{ext}}
=
\left\{
(y,p,\chi,\theta,\mathbf s)
\;\middle|\;
 p=\chi=\theta=0
\right\}.
\]
Après application de la regraine, tout état de ce régime prend la forme
\begin{equation}
\label{eq:annexe-H-etat-reduit}
X_{\mathrm r}(\mathbf s)
=
(y_0,0,0,0,\mathbf s).
\end{equation}
Deux histoires différentes qui aboutissent au même $\mathbf s$ donnent
donc le même état réduit au départ de la sonde. La conclusion
comportementale exige encore que toute variable interne omise soit
inerte sous la sonde, conformément à (E3) ci-dessous.

Le mot \emph{exacte} est essentiel. L'annulation numérique approchée des
résidus conduit à un problème de stabilité, traité seulement sous forme
de proposition conditionnelle à la fin de l'appendice.

\section{Hypothèses de réduction}
\label{annexe-H-hypotheses}

Le théorème utilise peu d'hypothèses, mais leur rôle doit être séparé.

\paragraph{O1 --- Déterminisme de l'exécution.}
Pour toute sonde admissible $\mathcal P$ et tout état de départ $X$, il
existe une exécution unique
\[
\operatorname{Exec}_{\mathcal P}(X).
\]
La lecture $\omega_\rho$ appliquée à cette exécution est donc elle aussi
univoque.

\paragraph{(E1) --- Regraine canonique.}
La regraine impose le même état mécanique $y_0$ à tous les systèmes
comparés et préserve les variables internes $(p,\chi,\theta,\mathbf s)$,
comme dans \eqref{eq:annexe-H-regraine}.

\paragraph{(E2) --- Extinction des variables non survivantes.}
Au départ de la sonde, les variables $(p,\chi,\theta)$ sont nulles.

\paragraph{(E3) --- Suffisance comportementale du résumé.}
Le résumé $(y,p,\chi,\theta,\mathbf s)$ est suffisant pour déterminer
l'évolution future sous la sonde déclarée : toute variable interne non
incluse dans ce résumé est soit fixée par le protocole, soit
comportementalement inerte sur la classe considérée.

L'hypothèse (E3) est la forme markovienne de la réduction. Elle est à
vérifier modèle par modèle : elle ne requiert pas que deux histoires
égalisées en $\mathbf s$ possèdent littéralement toutes les mêmes
variables internes, mais que les variables omises soient
comportementalement inertes sous la sonde déclarée. Dans le modèle
canonique, la seule variable additionnelle, $ssc$, peut différer entre
deux histoires ; sous la loi $\mathcal G_{\mathrm{lin}}$, elle
n'apparaît toutefois dans aucune équation active et est donc inerte,
comme établi dans la preuve de la Proposition~\ref{prop:I-etat-lent-aligne}. Sous une
loi lisant $ssc$, (E3) exigerait de l'inclure dans le résumé.

\begin{encadre}[title={Ce que le théorème ne suppose pas}]
Aucune des propriétés suivantes n'est requise :
\begin{itemize}
\item monotonie de $\Omega_{\mathcal P,\rho}$ ;
\item existence ou unicité d'une commutation ;
\item monotonie de l'écriture du préfixe ;
\item linéarité de la loi de mémoire ;
\item stabilité de la classe d'histoires par continuation.
\end{itemize}
Ces propriétés n'interviennent que dans certains corollaires.
\end{encadre}

\section{Réduction au contenu survivant}
\label{annexe-H-theoreme}

\resultstep{thm:H-contenu-survivant}
\begin{encadre}[title={Théorème \theproposition{} --- Réduction au contenu survivant}]
Soient $\mathcal P$ une sonde admissible et $\rho$ une résolution
d'observation. Sous O1 et (E1)--(E3), pour tout système dont l'état au moment
de la regraine appartient à $\Sigma_{\mathrm{ext}}$, il existe une
fonction
\[
\Omega_{\mathcal P,\rho}:\mathcal S\longrightarrow\mathcal Y_\rho
\]
telle que
\begin{equation}
\label{eq:annexe-H-theoreme-reduction}
\boxed{
\mathcal R_X^{\rho}(\mathcal P)
=
\Omega_{\mathcal P,\rho}(\mathbf s_{\mathrm r}),
}
\end{equation}
où $\mathbf s_{\mathrm r}$ est le contenu survivant à la regraine.
Pour une famille alignée $\mathbf s=s\mathbf u_1$, la fonction
$\Omega_{\mathcal P,\rho}$ se réduit à une fonction d'une variable
réelle.
\end{encadre}

\paragraph{Démonstration.}
Soit $X$ un système issu d'une histoire quelconque et supposons que son
état au moment de la regraine appartienne à
$\Sigma_{\mathrm{ext}}$. Par (E2), on a
\[
p=\chi=\theta=0.
\]
Par (E1), la regraine remplace la configuration mécanique par $y_0$ et
préserve le contenu survivant. L'état complet au départ de la sonde est
donc
\[
X_{\mathrm r}(\mathbf s_{\mathrm r})
=
(y_0,0,0,0,\mathbf s_{\mathrm r}).
\]
Par (E3), aucune information supplémentaire sur l'histoire n'est requise
pour déterminer la dynamique future. Par O1, l'exécution
\[
\operatorname{Exec}_{\mathcal P}
\bigl(X_{\mathrm r}(\mathbf s_{\mathrm r})\bigr)
\]
est une fonction univoque de $\mathbf s_{\mathrm r}$. En composant cette
application avec la lecture $\omega_\rho$, on définit
\begin{equation}
\label{eq:annexe-H-definition-omega}
\Omega_{\mathcal P,\rho}(\mathbf s)
:=
\omega_\rho\!\left(
\operatorname{Exec}_{\mathcal P}
(y_0,0,0,0,\mathbf s)
\right).
\end{equation}
Il vient immédiatement
\[
\mathcal R_X^{\rho}(\mathcal P)
=
\Omega_{\mathcal P,\rho}(\mathbf s_{\mathrm r}).
\]
Dans une famille alignée, $\mathbf s=s\mathbf u_1$ et la restriction de
$\Omega_{\mathcal P,\rho}$ à cette droite est une fonction du scalaire
$s$. \hfill$\square$

\subsection{Nature de la fonction \texorpdfstring{$\Omega$}{Omega}}
\label{annexe-H-nature-Omega}

La fonction $\Omega_{\mathcal P,\rho}$ n'est pas une quantité ajoutée au
modèle. Elle est la carte de réponse de la sonde nue, obtenue en
injectant directement un contenu dans un système canonique éteint puis
en appliquant la sonde :
\[
\mathbf s
\longmapsto
(y_0,0,0,0,\mathbf s)
\longmapsto
\operatorname{Exec}_{\mathcal P}
\longmapsto
\mathcal Y_\rho.
\]
Elle peut être calculée analytiquement, mesurée numériquement ou estimée
expérimentalement. Le théorème affirme que cette carte nue suffit ensuite
à prédire toute expérience préfixée qui rejoint le régime d'extinction.

Pour une lecture binaire d'appropriation, on peut écrire
\[
\Omega_{\mathcal P,\rho}(\mathbf s)
=
\mathbf 1_{A_{\mathcal P,\rho}}(\mathbf s),
\]
où
\[
A_{\mathcal P,\rho}
=
\left\{
\mathbf s\in\mathcal S
\;\middle|\;
\Omega_{\mathcal P,\rho}(\mathbf s)=1
\right\}
\]
est l'ensemble d'appropriation de la sonde sur la section éteinte. Sa
frontière
\[
\partial A_{\mathcal P,\rho}
\]
peut posséder plusieurs composantes. Le théorème~\ref{thm:H-contenu-survivant} ne la suppose ni
connexe ni réduite à un point.

\section{Préfixes éteints et carte totale d'écriture}
\label{annexe-H-prefixes}

\subsection{Écriture produite par un préfixe}
\label{annexe-H-ecriture}

Dans toute cette section, sauf décomposition explicite, $\mathcal A$
désigne le protocole préfixe complet : épisodes actifs, délai
d'extinction et regraine déclarée. Lorsque ces trois morceaux doivent
être distingués, on écrit $\mathcal A_{\mathrm{act}}\cdot\mathcal D\cdot\mathcal G$.

Considérons une famille d'états initiaux alignés
\[
X_0(s_0)=(y_{\mathrm{init}},p_{\mathrm{init}},
\chi_{\mathrm{init}},\theta_{\mathrm{init}},s_0\mathbf u_1)
\]
et un préfixe $\mathcal A$ suivi d'un délai d'extinction et d'une
regraine. Lorsque le protocole rejoint $\Sigma_{\mathrm{ext}}$, il
induit une application
\[
S_{\mathcal A}:I\longrightarrow\mathbb R,
\]
appelée \emph{carte totale d'écriture}, telle que
\begin{equation}
\label{eq:annexe-H-carte-ecriture}
\mathbf s_{\mathrm r}
=
S_{\mathcal A}(s_0)\,\mathbf u_1.
\end{equation}

Lorsque l'on souhaite isoler l'apport du préfixe, on écrit
\begin{equation}
\label{eq:annexe-H-delta}
S_{\mathcal A}(s_0)
=
s_0+\delta_{\mathcal A}(s_0).
\end{equation}
La quantité $\delta_{\mathcal A}$ est l'écriture nette relativement au
contenu initial. L'objet dynamique pertinent reste toutefois la carte
totale $S_{\mathcal A}$. En particulier, une propriété de monotonie de
$\delta_{\mathcal A}$ ne se confond pas avec une propriété de monotonie
de $S_{\mathcal A}$.

Cette distinction évite une erreur fréquente : le passage d'un motif
initial à un motif final dépend de
\[
\Omega_{\mathcal P,\rho}\circ S_{\mathcal A},
\]
et non de $\delta_{\mathcal A}$ considérée isolément.

\subsection{Formule de composition}
\label{annexe-H-formule-composition}

Le théorème~\ref{thm:H-contenu-survivant} donne immédiatement la formule suivante.

\resultstep{cor:H-prediction-prefixe}
\begin{encadre}[title={Corollaire \theproposition{} --- Prédiction d'une sonde après un préfixe éteint}]
Sous les hypothèses du Théorème~\ref{thm:H-contenu-survivant}, pour une lecture compatible avec la
segmentation déclarée au sens de l'Appendice~G, si le préfixe complet
$\mathcal A$ conduit au contenu
$S_{\mathcal A}(s_0)\mathbf u_1$, alors
\begin{equation}
\label{eq:annexe-H-composition}
\boxed{
\mathcal R_{X_0(s_0)}^{\rho}
(\mathcal A\cdot\mathcal P)
=
\Omega_{\mathcal P,\rho}
\bigl(S_{\mathcal A}(s_0)\bigr).
}
\end{equation}
Lorsque $S_{\mathcal A}(s_0)=s_0+\delta_{\mathcal A}(s_0)$,
\[
\mathcal R_{X_0(s_0)}^{\rho}
(\mathcal A\cdot\mathcal P)
=
\Omega_{\mathcal P,\rho}
\bigl(s_0+\delta_{\mathcal A}(s_0)\bigr).
\]
\end{encadre}

\paragraph{Démonstration.}
Après le préfixe, le délai et la regraine, l'état complet est
\[
(y_0,0,0,0,S_{\mathcal A}(s_0)\mathbf u_1).
\]
L'application du Théorème~\ref{thm:H-contenu-survivant} à cet état donne directement
\eqref{eq:annexe-H-composition}. \hfill$\square$

Cette formule est prédictive : une fois mesurées séparément la carte nue
$\Omega_{\mathcal P,\rho}$ et la carte d'écriture
$S_{\mathcal A}$, aucune nouvelle simulation du système préfixé n'est
nécessaire pour déterminer sa réponse.

\section{Cas à commutation unique et équation de seuil}
\label{annexe-H-seuil}

\subsection{Hypothèse locale de commutation unique}
\label{annexe-H-commutation-unique}

Supposons la lecture binaire et considérons un intervalle $J$ de
contenus parcouru par $S_{\mathcal A}$. On dit que la sonde possède une
commutation unique sur $J$ s'il existe $s^{\star}\in J$ tel que, à une
orientation près,
\begin{equation}
\label{eq:annexe-H-omega-seuil}
\Omega_{\mathcal P,\rho}(s)
=
\mathbf 1_{[s^{\star},+\infty)}(s)
\qquad
\text{pour }s\in J.
\end{equation}
Cette hypothèse est locale : la même sonde peut avoir d'autres
commutations hors de $J$.

\resultstep{cor:H-equation-seuil}
\begin{encadre}[title={Corollaire \theproposition{} --- Équation de seuil}]
Sous l'hypothèse de commutation unique
\eqref{eq:annexe-H-omega-seuil}, les points de bascule $b$ de la famille
préfixée sont les solutions de
\begin{equation}
\label{eq:annexe-H-seuil-general}
\boxed{
S_{\mathcal A}(b)=s^{\star},
\qquad
b+\delta_{\mathcal A}(b)=s^{\star}.
}
\end{equation}
Si, de plus, l'écriture est constante sur la famille,
$\delta_{\mathcal A}(s_0)=\delta$, alors $\delta$ est un invariant du
préfixe complet et l'équation prend sa forme historique
\begin{equation}
\label{eq:annexe-H-seuil-translation}
\boxed{
b+\delta=s^{\star}.
}
\end{equation}
\end{encadre}

\paragraph{Démonstration.}
Par le Corollaire~\ref{cor:H-prediction-prefixe}, la réponse préfixée vaut
$\Omega_{\mathcal P,\rho}(S_{\mathcal A}(s_0))$. Dans la plage où
\eqref{eq:annexe-H-omega-seuil} est valide, la valeur de la réponse
change exactement lorsque l'argument de $\Omega$ traverse
$s^{\star}$. Les points de bascule sont donc les antécédents de
$s^{\star}$ par $S_{\mathcal A}$. Le cas constant donne
\eqref{eq:annexe-H-seuil-translation}. \hfill$\square$

L'équation $b+\delta=s^{\star}$ n'est donc pas le résultat fondamental.
Elle est la forme particulière de
\[
\Omega\circ S
\]
lorsque la frontière de la sonde est ponctuelle sur la plage étudiée et
que l'écriture du préfixe est une translation.

\subsection{Nombre de seuils de la famille préfixée}
\label{annexe-H-nombre-seuils}

Même si $\Omega$ possède une unique commutation, l'équation
$S_{\mathcal A}(s_0)=s^{\star}$ peut avoir plusieurs solutions lorsque
$S_{\mathcal A}$ n'est pas injective. Réciproquement, même si
$S_{\mathcal A}$ est strictement monotone, plusieurs composantes de la
frontière de $A_{\mathcal P,\rho}$ peuvent engendrer plusieurs
commutations de la famille préfixée.

Ainsi, l'unicité du seuil final exige au moins deux contrôles distincts :
\begin{itemize}
\item la structure de la frontière de la sonde nue sur la plage atteinte ;
\item la géométrie de la carte totale d'écriture du préfixe.
\end{itemize}

\section{Proposition~\ref{prop:t3-compression-sous-critique} comme cas de compression}
\label{annexe-H-compression-sous-critique}

La Proposition~\ref{prop:t3-compression-sous-critique} concerne des histoires sous-critiques séparées par
des délais assez longs pour ramener les variables non mnésiques à leur
base. Dans cette classe, la dynamique du contenu est linéaire et les
résidus des épisodes se superposent avec leurs poids de récence.

Soit une histoire
\[
\mathcal H
=
E_1\cdot D_1\cdot E_2\cdots D_{n-1}\cdot E_n
\]
qui ne change pas de bassin, dont la regraine ramène les systèmes
comparés au même état autonome, et pour laquelle, avant la sonde,
\[
p=\chi=\theta=0.
\]
Supposons de plus que la loi d'évolution de $\mathbf s$ soit linéaire
sur cette classe et que l'épisode courant n'agisse sur la source du
contenu qu'à travers ses propres paramètres. Il existe alors des
résidus élémentaires $\mathbf r_i$ et des poids de récence $w_i$ tels
que
\begin{equation}
\label{eq:annexe-H-superposition}
\mathbf s_{\mathrm r}(\mathcal H)
=
\sum_{i=1}^{n}w_i\mathbf r_i.
\end{equation}

\resultstep{cor:H-compression-sous-critique}
\begin{encadre}[title={Corollaire \theproposition{} --- Compression sous-critique, Proposition~\ref{prop:t3-compression-sous-critique}}]
Dans la classe précédente, deux histoires $\mathcal H$ et
$\mathcal H'$ qui produisent le même contenu survivant vérifient, pour
toute sonde admissible $\mathcal P$ placée après extinction et regraine,
\begin{equation}
\label{eq:annexe-H-compression-equivalence}
\mathbf s_{\mathrm r}(\mathcal H)
=
\mathbf s_{\mathrm r}(\mathcal H')
\quad\Longrightarrow\quad
\mathcal R^{\rho}_{X_0}
(\mathcal H\cdot\mathcal P)
=
\mathcal R^{\rho}_{X_0}
(\mathcal H'\cdot\mathcal P).
\end{equation}
En particulier, injecter le contenu
$\mathbf s_{\mathrm r}(\mathcal H)$ dans un système canonique éteint
reproduit la réponse de l'histoire complète.
\end{encadre}

\paragraph{Démonstration.}
Les hypothèses de sous-criticalité et de séparation assurent que les
deux histoires rejoignent le régime d'extinction. Par hypothèse, la
regraine les ramène au même état autonome et leurs contenus survivants sont
égaux. Leurs états complets au départ de la sonde coïncident donc. Le
Théorème~\ref{thm:H-contenu-survivant} donne l'égalité des réponses. La formule
\eqref{eq:annexe-H-superposition} fournit en outre un calcul explicite du
résumé, mais elle n'est pas nécessaire à l'implication
\eqref{eq:annexe-H-compression-equivalence}. \hfill$\square$

Ce corollaire précise la portée de \ref{prop:t3-compression-sous-critique}. La linéarité explique comment
calculer le contenu survivant ; la suffisance comportementale de ce
contenu vient du Théorème~\ref{thm:H-contenu-survivant}. Dès qu'un épisode change de bassin, que
les délais laissent survivre une phase ou une pression, ou qu'une autre
variable interne est active, la réduction à $\mathbf s$ seul n'est plus
garantie.

\section{Structure générale des motifs}
\label{annexe-H-motifs}

\subsection{Décomposition des sources de non-monotonie}
\label{annexe-H-nonmonotonie}

Dans la famille scalaire, la réponse préfixée est exactement
\[
f(s_0)
=
\Omega_{\mathcal P,\rho}
\bigl(S_{\mathcal A}(s_0)\bigr).
\]
Toute alternance de la réponse est donc produite par l'intersection de
la courbe $S_{\mathcal A}(I)$ avec la frontière de l'ensemble
$A_{\mathcal P,\rho}$.

Pour une lecture binaire, deux mécanismes suffisent à engendrer des
motifs complexes :
\begin{enumerate}
\item la frontière de la sonde nue possède plusieurs composantes sur la
  plage atteinte ;
\item la carte totale d'écriture revient plusieurs fois sur une même
  composante de frontière, par exemple parce qu'elle n'est pas
  monotone.
\end{enumerate}

Cette dichotomie doit être comprise au niveau de la composition
$\Omega\circ S$. Elle ne permet pas d'attribuer un motif à la seule
fonction $\delta$. Une écriture nette $\delta$ croissante peut produire
une carte totale $S=s_0+\delta(s_0)$ croissante, et la réponse peut
malgré tout être anti-monotone si $\Omega$ possède une île de repos ou
d'appropriation. Inversement, une sonde nue monotone peut donner
plusieurs bascules si $S$ traverse plusieurs fois son seuil.

\resultstep{cor:H-localisation-bascules}
\begin{encadre}[title={Corollaire \theproposition{} --- Localisation exhaustive des bascules}]
Sous les hypothèses du Théorème~\ref{thm:H-contenu-survivant}, l'ensemble des points de bascule
de la famille préfixée est contenu dans
\begin{equation}
\label{eq:annexe-H-preimage-frontiere}
S_{\mathcal A}^{-1}
\bigl(\partial A_{\mathcal P,\rho}\bigr).
\end{equation}
Si $\Omega_{\mathcal P,\rho}$ est localement constante sur chaque
composante de
$\mathcal S\setminus\partial A_{\mathcal P,\rho}$ et si
$S_{\mathcal A}$ est continue, les changements effectifs de réponse ne
peuvent se produire qu'en ces antécédents.
\end{encadre}

\paragraph{Démonstration.}
Si $S_{\mathcal A}(s_0)$ appartient à l'intérieur d'une composante sur
laquelle $\Omega$ est constante, la continuité de $S_{\mathcal A}$
garantit que, dans un voisinage assez petit de $s_0$, son image reste
dans cette composante ; la réponse ne change pas. Toute bascule doit
donc correspondre à un point dont l'image rencontre la frontière.
\hfill$\square$

L'inclusion peut être stricte : une tangence à la frontière peut ne pas
changer la valeur binaire, et certaines composantes de frontière peuvent
être inaccessibles à la carte d'écriture.

\subsection{Îles de la sonde nue}
\label{annexe-H-iles}

Une \emph{île d'appropriation} est une composante bornée de
$A_{\mathcal P,\rho}$ sur l'axe ou la section étudiée. Une \emph{île de
repos} est une composante bornée de son complément. Ces objets
appartiennent déjà à la géométrie de la sonde nue. Ils ne sont donc pas,
en général, des artefacts produits par la composition avec une amorce.

Le rôle du préfixe est de déterminer le chemin suivi dans l'espace des
contenus. Une carte d'écriture monotone peut traverser successivement
une entrée puis une sortie d'île et produire un motif apparemment
anti-monotone, sans aucune inversion de l'écriture.

\section{Extinction approchée}
\label{annexe-H-extinction-approchee}

Le Théorème~\ref{thm:H-contenu-survivant} est exact sur $\Sigma_{\mathrm{ext}}$. En pratique, les
variables $(p,\chi,\theta)$ peuvent n'être que petites. On note
\[
z=(p,\chi,\theta,\mathbf s),
\qquad
z_0=(0,0,0,\mathbf s),
\]
et l'on suppose
\[
\|(p,\chi,\theta)\|\leq\varepsilon.
\]

La proximité des états initiaux n'implique pas à elle seule l'égalité
des lectures binaires : près d'une frontière de commutation, une
perturbation arbitrairement petite peut changer l'issue. Une formulation
correcte doit donc exclure un voisinage de la frontière et contrôler la
continuité de l'exécution et de la lecture avant seuillage.

\resultstep{prop:H-extinction-approchee}
\begin{encadre}[title={Proposition \theproposition{} --- Extinction approchée, énoncé conditionnel}]
Supposons que, sur un compact $K$ de contenus, l'application qui associe
à l'état initial une observable réelle pré-binaire soit continue, et que
la frontière de commutation soit traversée transversalement hors des
points de changement de calendrier. Alors il existe une fonction
$\eta(\varepsilon)\to0$ telle que, si
\[
\|(p,\chi,\theta)\|\leq\varepsilon,
\]
la réponse du système coïncide avec
$\Omega_{\mathcal P,\rho}(\mathbf s)$ pour tout contenu situé à une
distance supérieure à $\eta(\varepsilon)$ de la frontière pertinente.
Aucune borne uniforme n'est revendiquée au voisinage des changements de
calendrier de la sonde.
\end{encadre}

Cette proposition fixe la forme attendue du résultat, mais sa preuve
quantitative dépend de bornes de stabilité propres au modèle et reste
rattachée à QO-74. Elle ne doit pas être utilisée comme théorème acquis
dans le corps du manuscrit.

\section{Vérifications numériques associées}
\label{annexe-H-verifications}

Les calculs du Tome III ne constituent pas la preuve du Théorème~\ref{thm:H-contenu-survivant},
mais ils vérifient ses conséquences sur le modèle de référence.

\subsection{Mesure directe de la sonde nue}
\label{annexe-H-verification-omega}

Pour la sonde de durée $T_2=80$ et la résolution binaire de bassin, le
scan direct de la famille canonique
\[
(y_0,0,0,0,s\mathbf u_1)
\]
retrouve une commutation à
\[
s_{80}^{\star}\simeq0{,}01627
\]
sur la plage étudiée. Cette valeur est obtenue sans préfixe. Elle
coïncide avec celle qui avait d'abord été inférée à partir de plusieurs
histoires amorcées, ce qui fournit un contrôle indépendant de la
réduction.

\subsection{Prédiction par composition}
\label{annexe-H-verification-prediction}

Pour deux préfixes éteints, la carte d'écriture
$S_{\mathcal A}(s_0)$ a été mesurée point par point, puis injectée dans
la carte nue $\Omega_{80,\rho}$. La prédiction
\[
s_0
\longmapsto
\Omega_{80,\rho}
\bigl(S_{\mathcal A}(s_0)\bigr)
\]
coïncide cellule par cellule avec la carte obtenue en simulant
l'histoire complète. Cette vérification distingue la réduction d'une
simple reformulation a posteriori : les deux constituants sont mesurés
séparément puis recomposés.

\subsection{Commutations multiples}
\label{annexe-H-verification-iles}

La variation de la durée de sonde montre que la carte nue peut elle-même
posséder plusieurs composantes. Dans les calculs de référence, une sonde
présente une île d'appropriation et une autre une île de repos interne.
La réponse non monotone observée après un préfixe s'explique alors par la
traversée de cette île par une carte totale d'écriture croissante. Le
prédicat pointwise
\[
\Omega_{\mathcal P,\rho}
\bigl(S_{\mathcal A}(s_0)\bigr)
\]
reproduit toutes les cellules du témoin, tandis qu'un prédicat à seuil
unique échoue sur une partie d'entre elles.

Les valeurs numériques et scripts de reproductibilité relèvent de
l'annexe consacrée aux méthodes. Ici, leur fonction est seulement de
montrer que les deux mécanismes abstraits de la section précédente sont
effectivement réalisés par le modèle.

\section{Limites de portée}
\label{annexe-H-limites}

Le Théorème~\ref{thm:H-contenu-survivant} ne s'applique pas dès que l'état au départ de la sonde
n'est plus déterminé par le seul contenu survivant. Les cas suivants
sortent notamment de son domaine direct.

\paragraph{Délais courts.}
Une phase $\theta$, une pression $p$ ou une charge d'interface $\chi$
peut rester active. Le résumé pertinent devient alors au moins
$(p,\chi,\theta,\mathbf s)$.

\paragraph{Histoires conditionnelles.}
Un épisode peut changer de bassin. La regraine mécanique ne ramène pas
nécessairement les systèmes à un même état autonome, et le bassin atteint
doit être inclus dans le résumé.

\paragraph{Plasticité de la loi.}
Si l'histoire modifie une loi d'action ou une variable plastique
supplémentaire $\gamma$, l'hypothèse (E3) est fausse pour l'état réduit
utilisé ici. La réduction peut être rétablie dans un état augmenté, mais
elle ne porte plus sur $\mathbf s$ seul.

\paragraph{Lectures historiques intégrales.}
Une lecture qui compte des événements produits pendant le préfixe ne se
réduit pas à la réponse de la sonde nue. Le résultat porte sur la lecture
de la sonde après regraine, ou sur une lecture terminale compatible avec
la segmentation déclarée, conformément à l'Appendice~G.

\paragraph{Extinction seulement numérique.}
Une petite norme résiduelle ne garantit pas une même issue au voisinage
d'une frontière. La proposition~\ref{prop:H-extinction-approchee} exige des hypothèses de
stabilité supplémentaires.

\section{Dépendances logiques et usages}
\label{annexe-H-dependances}

Le rôle des résultats peut être résumé ainsi :
\[
\begin{array}{c}
\text{regraine canonique}
+
\text{extinction exacte}
+
\text{fermeture de l'état}
+
\text{déterminisme}
\\[2mm]
\Downarrow
\\[2mm]
\mathcal R(\mathcal P)=\Omega_{\mathcal P,\rho}(\mathbf s_{\mathrm r})
\qquad\text{(Théorème~\ref{thm:H-contenu-survivant})}
\\[2mm]
\Downarrow
\\[2mm]
\mathcal R(\mathcal A\cdot\mathcal P)
=
\Omega_{\mathcal P,\rho}\circ S_{\mathcal A}
\\[2mm]
\Downarrow
\\[2mm]
\begin{cases}
S_{\mathcal A}(b)=s^{\star},
&\text{si la sonde a une commutation unique},\\
\text{compression par }\mathbf s,
&\text{pour les histoires sous-critiques séparées},\\
S_{\mathcal A}^{-1}(\partial A_{\mathcal P,\rho}),
&\text{pour localiser les bascules générales}.
\end{cases}
\end{array}
\]

Le Théorème~\ref{thm:H-contenu-survivant} est donc plus général que l'équation de seuil et plus
fondamental que la superposition linéaire de la mémoire. Il sépare trois
questions qui avaient initialement tendance à se confondre :
\begin{enumerate}
\item quel état survit au préfixe ;
\item comment le préfixe écrit cet état ;
\item comment la sonde répond à cet état.
\end{enumerate}
La première relève du régime d'extinction, la deuxième de la carte
$S_{\mathcal A}$, la troisième de
$\Omega_{\mathcal P,\rho}$.

\section{Table des notations}
\label{annexe-H-notations}

\begin{longtable}{>{\raggedright\arraybackslash}p{0.25\textwidth}
                  >{\raggedright\arraybackslash}p{0.67\textwidth}}
\toprule
Notation & Signification \\
\midrule
\endfirsthead
\toprule
Notation & Signification \\
\midrule
\endhead
$X=(y,p,\chi,\theta,\mathbf s)$
& état complet utilisé dans la réduction \\
$y_0$
& configuration mécanique canonique imposée par la regraine \\
$\mathcal G$
& épisode de regraine, appartenant au protocole \\
$\Sigma_{\mathrm{ext}}$
& régime d'extinction exacte $p=\chi=\theta=0$ \\
$\mathbf s_{\mathrm r}$
& contenu survivant au départ de la sonde \\
$\mathcal P$
& sonde appliquée après extinction et regraine \\
$\Omega_{\mathcal P,\rho}$
& carte de réponse de la sonde nue à la résolution $\rho$ \\
$A_{\mathcal P,\rho}$
& ensemble des contenus donnant l'appropriation pour la sonde nue \\
$\partial A_{\mathcal P,\rho}$
& frontière de commutation sur l'espace des contenus \\
$\mathcal A$
& préfixe d'histoire étudié avant la sonde \\
$S_{\mathcal A}$
& carte totale d'écriture produite par le préfixe \\
$\delta_{\mathcal A}$
& écriture nette définie par $S_{\mathcal A}(s_0)=s_0+\delta_{\mathcal A}(s_0)$ \\
$s^{\star}$
& commutation de la sonde nue sur une plage où elle est unique \\
$\eta(\varepsilon)$
& largeur d'exclusion attendue autour de la frontière en extinction approchée \\
\bottomrule
\end{longtable}

\section{Conclusion}
\label{annexe-H-conclusion}

La réduction au contenu survivant repose sur une égalité d'états, non
sur une approximation phénoménologique. Une fois les variables
non survivantes éteintes et la configuration mécanique regrainée, deux
histoires qui laissent le même contenu donnent le même état complet au
départ de la sonde ; le déterminisme impose alors la même réponse.

La mécanique générale se résume à la composition
\[
\boxed{
\text{réponse préfixée}
=
\text{carte de la sonde nue}
\circ
\text{carte totale d'écriture du préfixe}.
}
\]
Les seuils uniques, les translations, les îles et les motifs
anti-monotones sont différentes géométries de cette même composition.
Le résultat prépare ainsi l'étape suivante : remplacer la section
éteinte, paramétrée par $\mathbf s$, par l'état lent complet
$(p,\chi,\theta,\mathbf s)$ et étudier la frontière d'appropriation qui
s'y déploie.
